% TOP_PROBLEM: {"id": 2823, "title": "Kirby Problem 3.25", "category_id": 7, "category": {"id": 7, "name": "topology", "display_name": "Topology", "description": "Properties preserved under continuous deformations.", "slug": "topology", "order_index": 7, "created_at": "2026-07-31T15:26:25.671Z"}, "status": "open", "problem_number": "KP-3.25", "set_id": 11}
% FIRST_POSTED: 2026-10-01
% ENTRY_KIND: statement_only
% UnsolvedMath ID 2823; source catalogue code KP-3.25
% !TeX program = lualatex
% Definition and sources only; this file is not an LLM solution attempt.
\documentclass[11pt,a4paper]{article}
\usepackage[margin=25mm]{geometry}
\usepackage{fontspec}
\setmainfont{lmroman10-regular.otf}[BoldFont=lmroman10-bold.otf,
 ItalicFont=lmroman10-italic.otf,BoldItalicFont=lmroman10-bolditalic.otf]
\usepackage{amsmath,amssymb,mathtools}
\usepackage{unicode-math}
\setmathfont{latinmodern-math.otf}
\AtBeginDocument{\let\setminus\smallsetminus}
\usepackage{xurl,hyperref}
\hypersetup{colorlinks=true,urlcolor=blue}
\setcounter{secnumdepth}{0}
\setlength{\parindent}{0pt}
\setlength{\parskip}{5pt}
\setlength{\emergencystretch}{3em}
\begin{document}
\section*{Kirby Problem 3.25}\label{2823}
{\small UnsolvedMath ID 2823 · KP-3.25 · Topology · Kirby's Problems in Low-Dimensional Topology}
\subsection{Definitions and mathematical statement}
Let $S$ be an oriented finite-type surface with $\chi(S)<0$. A marked
quasi-Fuchsian structure is a discrete faithful representation
$\rho:\pi_1(S)\to\operatorname{PSL}_2(\mathbb C)$ obtained by a
quasiconformal deformation of a Fuchsian representation, modulo conjugacy;
puncture loops remain parabolic. Its manifold $M_\rho$ is homeomorphic to
$S\times\mathbb R$. The convex core is the quotient of the convex hull
of the limit set in $\partial_\infty\mathbb H^3$.

The two marked boundary components of this core have intrinsic complete
hyperbolic metrics $m_+,m_-$. They are pleated along measured geodesic
laminations $\lambda_+,\lambda_-$, where a lamination is a closed union
of disjoint complete geodesics and its transverse measure records the
exterior bending angle. Thus one obtains maps
\[
 \mu:\mathcal QF(S)\to\mathcal T(S)^2,\qquad
 \beta:\mathcal QF(S)\setminus\mathcal F(S)\to\mathcal{ML}(S)^2.
\]
Here $\mathcal T(S)$ is marked hyperbolic-metric space, $\mathcal{ML}(S)$
is measured-lamination space, and $\mathcal F(S)$ is the Fuchsian locus.

(a) Is $\mu$ injective? (b) When $S$ has punctures, is $\beta$ injective?
Equality includes the markings and the ordering of the two ends. In (a)
a Fuchsian core is read as two copies of its intrinsic surface metric.
In (b) the Fuchsian locus is excluded because its bending measures vanish.
The unpunctured bending-uniqueness case is not being reasserted as open.
\subsection{Short English statement}
Do the two intrinsic convex-core metrics determine a marked quasi-Fuchsian manifold, and do the two bending measures determine it in the punctured non-Fuchsian case?
\subsection{Sources}
\begin{itemize}
\item[S1] ULAM.AI and the UnsolvedMath Contributors, supplied problems.json, record 2823 (KP-3.25), Kirby's Problems in Low-Dimensional Topology. Catalogue identity and source transcription; CC BY 4.0.
\url{https://huggingface.co/datasets/ulamai/UnsolvedMath}.
\item[S2] K3: A New Problem List in Low-Dimensional Topology, Problem 3.25. Expanded formulation of the supplied catalogue statement.
\url{https://math.berkeley.edu/sites/default/files/surv-295-ruberman-watermarked-author-pdf.pdf}.
\end{itemize}
\end{document}
